Parmi les méthodes d'amas quantiques dédiées à la résolution des modèles de type Hubbard, la CDMFT couplée à un solutionneur d'impureté par diagonalisation exacte (ED-CDMFT) offre un accès privilégié aux propriétés spectrales à température nulle. Or, la dimension de l'espace de Hilbert croît comme $4^{N_{\rm c} + N_{\rm b}}$, rendant rapidement inaccessibles les systèmes nécessitant un bain très détaillé pour bien représenter l'environnement. C'est précisément à cette difficulté que la SB-CDMFT, nouvelle méthode présentée dans ce travail, prétend à s'attaquer. L'idée centrale repose sur la subdivision du bain discret de $N_{\rm b}$ orbitales en $N_{\rm sb}$ sous-bains indépendants. Chaque sous-bain est couplé à l'amas et les $N_{\rm sb}$ problèmes d'impureté résultants sont résolus de façon indépendante. La self-énergie du réseau est finalement approximée comme la moyenne des contributions de ces différents problèmes d'impureté. Pour chaque diagonalisation, le facteur de réduction $\mathcal{R}$ de la taille de l'espace de Hilbert est exponentiel~:
\begin{align}
    \mathcal{R} = \frac{4^{N_{\rm b}(N_{\rm sb} - 1)/N_{\rm sb}}}{N_{\rm sb}}.
\end{align}

Plusieurs idées émergent naturellement de ce travail. Ce concept de sous-bains ouvre la possibilité d'un solutionneur d'impureté par diagonalisation exacte pour la DCA, habituellement limitée par le problème de signe fermionique présent dans les solutionneurs d'impureté basés sur le Monte-Carlo. L'application de la SB-CDMFT à des matériaux réalistes comme les cuprates représente un objectif à plus long terme, pour lequel le gain calculatoire de la méthode pourrait s'avérer intéressant~\cite{bbl-prx}.
